% !TeX program = lualatex
\documentclass[11pt, a4paper]{article}

% --- EVRENSEL ÖNSÖZ BLOĞU ---
\usepackage[a4paper, top=2.5cm, bottom=2.5cm, left=2cm, right=2cm]{geometry}
\usepackage{fontspec}
\usepackage[turkish, bidi=basic, provide=*]{babel}

% Noto Sans font ayarları
\babelprovide[import, onchar=ids fonts]{turkish}
\babelprovide[import, onchar=ids fonts]{english}
\babelfont{rm}{Noto Sans}

% Liste simgeleri için
\usepackage{enumitem}
\setlist[itemize]{label=-}

% Matematik ve Tablo Paketleri
\usepackage{amsmath, amssymb, amsthm}
\usepackage{booktabs}
\usepackage{array}

% Bölüm numaralandırmasını kaldırmak için (isteğe bağlı)
\setcounter{secnumdepth}{0}

\begin{document}

\title{Depo Toplama Rotası Optimizasyonu: Matematiksel Model}
\author{}
\date{}
\maketitle

\section{Kümeler ve İndisler}

\begin{itemize}
    \item $i \in I$: Ürün artikel kodları kümesi, $I = \{1, \dots, 3388\}$
    \item $j_1 \in J_1$: Katlar, $J_1 = \{1, 2, 3, 4, 5, 6\}$
    \item $j_2 \in J_2$: Koridorlar, $J_2 = \{1, \dots, 27\}$
    \item $j_3 \in J_3$: Koridor tarafı (Sol: 0, Sağ: 1), $J_3 = \{0, 1\}$
    \item $j_4 \in J_4$: Sütunlar, $J_4 = \{1, \dots, 20\}$
    \item $j_5 \in J_5$: Raflar, $J_5 = \{1, 2, 3, 4, 5, 6, 7\}$
    \item $j_6 \in J_6$: THM (Taşıma Hazırlama Merkezi) kodu
    \item $j = (j_1, j_2, j_3, j_4, j_5, j_6)$: 6 boyutlu tam lokasyon indisi, $j \in J$
    \item $n, m = (j_1, j_2, j_4)$: 3 boyutlu fiziksel konum indisi, $n, m \in N$
    \item $N(j_1)$: $j_1$ katındaki fiziksel konumlar kümesi
    \item $0_{j_1}$: $j_1$ katı için sanal giriş/çıkış deposu
    \item $A(j_1)$: $j_1$ katındaki konumlar ve $0_{j_1}$ deposu arasındaki yönlü rota yayları kümesi
\end{itemize}

\section{Parametreler}

\begin{itemize}
    \item $d_{n,m}$: Aynı kattaki $n$ konumu ile $m$ konumu arasındaki Manhattan mesafesi
    \item $d_{0,n}$, $d_{n,0}$: Kat deposu ile $n$ konumu arasındaki giriş/çıkış mesafesi
    \item $s_{ij}$: $i$ ürününün $j$ lokasyonundaki mevcut stok miktarı
    \item $t_i$: $i$ ürününden toplanması gereken toplam miktar
    \item $w_1, w_2, w_3$: Amaç fonksiyonu ağırlıkları (Sırasıyla mesafe, kutu sayısı ve kat sayısı için)
\end{itemize}

\section{Karar Değişkenleri}

\begin{itemize}
    \item $y_{ij} \in \{0\} \cup \mathbb{Z}_{\geq 1}$: $i$ ürününün $j$ lokasyonundan toplanan miktarı
    \item $u_n \in \{0, 1\}$: $n$ konumunun ziyaret edilip edilmeme durumu
    \item $v_{n,m} \in \{0, 1\}$: $(n,m) \in A(j_1)$ yayı üzerinden doğrudan gidilme durumu
    \item $p_n$: $n$ konumunun tur içindeki ziyaret sırası (MTZ kısıtı için)
    \item $z_{j_1} \in \{0, 1\}$: $j_1$ numaralı kata girilip girilmediği durumu
    \item $b_{j_6} \in \{0, 1\}$: $j_6$ numaralı THM kutusunun açılıp açılmadığı durumu
\end{itemize}

\section{Amaç Fonksiyonu}

\begin{equation}
    \min Z = w_1 \sum_{j_1 \in J_1} \sum_{(n,m) \in A(j_1)} d_{n,m} v_{n,m}
    + w_2 \sum_{j_6 \in J_6} b_{j_6}
    + w_3 \sum_{j_1 \in J_1} z_{j_1}
\end{equation}

\noindent \textbf{Terim Açıklamaları:}
\begin{itemize}
    \item \textbf{Terim 1:} Toplam yürünen Manhattan mesafesi.
    \item \textbf{Terim 2:} Toplam kullanılan (açılan) THM kutusu sayısı.
    \item \textbf{Terim 3:} Toplam girilen (aktifleşen) kat sayısı.
\end{itemize}

\section{Kısıtlar}

\begin{flalign}
    & 0 \leq y_{ij} \leq s_{ij} \cdot b_{j_6} & \forall i \in I, \forall j \in J \label{cons:stok} \\
    & y_{ij} \leq s_{ij} \cdot u_n & \forall i \in I, \forall n \in N, \forall j \in J(n) \label{cons:ziyaret} \\
    & \sum_{j \in J} y_{ij} = t_i & \forall i \in I \label{cons:talep} \\
    & u_n \leq \sum_{i \in I} \sum_{j \in J(n)} y_{ij} & \forall n \in N \label{cons:ziyaretters} \\
    & u_n \leq z_{j_1} & \forall j_1 \in J_1, \forall n \in N(j_1) \label{cons:kat} \\
    & z_{j_1} \leq \sum_{n \in N(j_1)} u_n & \forall j_1 \in J_1 \label{cons:katter} \\
    & b_{j_6} \leq \sum_{i \in I} \sum_{j \in J(j_6)} y_{ij} & \forall j_6 \in J_6 \label{cons:thmters} \\
    & \sum_{m:(n,m) \in A(j_1)} v_{n,m} = u_n & \forall j_1 \in J_1, \forall n \in N(j_1) \label{cons:flow1} \\
    & \sum_{m:(m,n) \in A(j_1)} v_{m,n} = u_n & \forall j_1 \in J_1, \forall n \in N(j_1) \label{cons:flow2} \\
    & \sum_{n \in N(j_1)} v_{0_{j_1},n} = z_{j_1} & \forall j_1 \in J_1 \label{cons:depotout} \\
    & \sum_{n \in N(j_1)} v_{n,0_{j_1}} = z_{j_1} & \forall j_1 \in J_1 \label{cons:depotin} \\
    & u_n \leq p_n \leq |N(j_1)| \cdot u_n & \forall j_1 \in J_1, \forall n \in N(j_1) \label{cons:mtzbound} \\
    & p_n - p_m + |N(j_1)| \cdot v_{n,m} \leq |N(j_1)| - 1 & \forall j_1 \in J_1, \forall n, m \in N(j_1), n \neq m \label{cons:mtz}
\end{flalign}

\noindent \textbf{Kısıt Açıklamaları:}
\begin{itemize}
    \item \textbf{(\ref{cons:stok})}: Pozitif toplama yapılan bir lokasyonun ilgili THM kutusu açılmış sayılır ve toplanan miktar stok miktarını geçemez.
    \item \textbf{(\ref{cons:talep})}: Her ürün için toplanması gereken toplam miktar sağlanmalıdır.
    \item \textbf{(\ref{cons:ziyaret})}: Ürün toplama yapılan her lokasyonun bulunduğu fiziksel konuma uğranmış olmalıdır.
    \item \textbf{(\ref{cons:ziyaretters})}: Ziyaret değişkeni yalnızca ilgili fiziksel konumda gerçekten toplama yapılıyorsa açılabilir.
    \item \textbf{(\ref{cons:kat})-(\ref{cons:katter})}: Kat değişkeni yalnızca o katta en az bir fiziksel konum ziyaret ediliyorsa açılabilir.
    \item \textbf{(\ref{cons:thmters})}: THM değişkeni yalnızca o THM altında gerçekten toplama yapılıyorsa açılabilir.
    \item \textbf{(\ref{cons:flow1})-(\ref{cons:flow2})}: Akış korunumu kısıtları. Bir noktaya girildiyse oradan çıkılmalıdır.
    \item \textbf{(\ref{cons:depotout})-(\ref{cons:depotin})}: Her aktif kat rotası sanal kat deposundan bir kez çıkar ve aynı depoya bir kez döner.
    \item \textbf{(\ref{cons:mtzbound})-(\ref{cons:mtz})}: Miller-Tucker-Zemlin (MTZ) alt tur eleme kısıtı her kat için bağımsız uygulanır.
\end{itemize}

\section{Uygulama Notu}

Kod tarafında her aktif kat için ayrı bir depo-bağlantılı rota kurulmaktadır. Bu sayede kat içi yürüyüş mesafesi, THM açma maliyeti ve aktif kat maliyeti ayrı ayrı ölçülebilir. $x_{ij}$ ikili değişkeni ayrıca tutulmaz; semi-integer $y_{ij}$ değişkeni $0$ veya pozitif tam sayı değer alarak aynı seçimi daha az değişkenle temsil eder.

\end{document}
